\documentclass[11pt]{article}
\usepackage[T1]{fontenc}
\usepackage[utf8]{inputenc}
\usepackage{lmodern}
\usepackage[margin=1in]{geometry}
\usepackage{microtype}
\emergencystretch=2em
\tolerance=1500
\hbadness=10000

\usepackage{amsmath,amssymb,amsthm,mathtools}
\usepackage{enumitem}
\usepackage{hyperref}
\usepackage{xurl}
\usepackage{tikz-cd}

\theoremstyle{plain}
\newtheorem{theorem}{Theorem}[section]
\newtheorem{lemma}[theorem]{Lemma}
\newtheorem{proposition}[theorem]{Proposition}
\newtheorem{corollary}[theorem]{Corollary}

\theoremstyle{definition}
\newtheorem{definition}[theorem]{Definition}
\newtheorem{remark}[theorem]{Remark}

\newcommand{\RS}{\mathrm{R\text{-}S}}
\newcommand{\cS}{\mathcal{S}}
\newcommand{\cF}{\mathcal{F}}
\newcommand{\cM}{\mathcal{M}}
\newcommand{\cC}{\mathcal{C}}
\newcommand{\cD}{\mathcal{D}}
\newcommand{\cU}{\mathcal{U}}
\newcommand{\cW}{\mathcal{W}}
\newcommand{\fM}{\mathfrak{M}}
\newcommand{\Meta}{\mathfrak{Meta}}
\newcommand{\Sint}{\cS_{\mathrm{int}}}
\newcommand{\Ob}{\mathrm{Ob}}
\newcommand{\Mor}{\mathrm{Mor}}
\newcommand{\Hom}{\mathrm{Hom}}
\newcommand{\Fun}{\mathrm{Fun}}
\newcommand{\Cl}{\mathrm{Cl}}
\newcommand{\Set}{\mathrm{Set}}
\newcommand{\Cat}{\mathrm{Cat}}
\newcommand{\Syn}{\mathrm{Syn}}
\newcommand{\Mod}{\mathrm{Mod}}
\newcommand{\Con}{\mathrm{Con}}
\newcommand{\ZFC}{\mathsf{ZFC}}
\newcommand{\MLTT}{\mathsf{MLTT}}
\newcommand{\HoTT}{\mathsf{HoTT}}
\newcommand{\ETT}{\mathsf{ETT}}
\newcommand{\II}{\mathbf{I}}
\DeclareMathOperator{\colim}{colim}
\DeclareMathOperator{\image}{image}
\DeclareMathOperator{\id}{id}
\DeclareMathOperator{\ev}{ev}

\title{The Absolute Foundation No-Go Theorem}

\author{%
  \footnotesize Max Sereda \\[-3pt]
  \scriptsize\textit{Independent researcher} \\[-3pt]
  \scriptsize\texttt{\href{mailto:max@gst.st}{max@gst.st}}%
}

\date{\scriptsize\today}

\begin{document}
\maketitle

\begin{abstract}
We state a relative boundary lemma for formally represented structures in a specified metatheory: if a candidate belongs to the selected derived class $\cS_S$, its identity is an interpretation into that class, so it cannot satisfy a condition forbidding every such interpretation. This is the precise content of the AFN-T obstruction used here; it does not exclude new foundations or operational refinements. We distinguish it from conditional classifier-comparison and slice constructions. Categoricity from the maximality axioms, strict intensional refinement, unbounded meta-stabilization, and an exact consistency-strength calculation remain open construction obligations. This summary follows the corrected main MSFS paper.

\medskip
\noindent\textbf{Keywords:} no-go theorems, categorical logic, $(\infty, n)$-categories, Lambek--Scott semantics, meta-classification.

\noindent\textbf{MSC2020:} 03A05, 03B30, 18A05, 18N10, 18N50, 18N60, 18N65.
\end{abstract}

\tableofcontents

\section{Introduction}\label{sec:intro}

Classical results of Cantor, Russell, G\"odel, Tarski, and Lawvere have different hypotheses and conclusions. The relative boundary lemma below is an identity obstruction for an explicitly defined comparison class; it does not reprove those results or subsume their specific diagonals.

\paragraph{Main result.} A specified presentation in $\cS_S$ reduces to itself by the identity. Thus the condition $\Pi_4$, which forbids all such reductions, fails (Theorem~\ref{thm:afnt}). No full faithfulness of arbitrary syntax-to-model evaluation is claimed.

\paragraph{Scope.} The same argument applies at each constructed categorical or metatheoretic stage where its hypotheses hold. It does not construct coherent higher-depth models or an unbounded tower of universes.

\paragraph{Comparison and refinement.} Theorems~\ref{thm:meta-cat} and~\ref{thm:meta-mult} give conditional comparison criteria. Theorem~\ref{thm:intensional} supplies an ordinary slice construction, whose equivalence-only display filtration is trivial. Further intensional data require their own construction and preservation proofs.

\paragraph{Open obligations.} Remarks~\ref{thm:stab}, \ref{thm:reflective}, and~\ref{thm:meta-stab} record missing stabilization and strength arguments. No exact equiconsistency with two inaccessible cardinals follows from choosing that ambient base. No machine-checked coverage of the revised claims is asserted without an exact statement mapping and a fresh replay of the companion corpus.

\paragraph{Conventions.} We work throughout in $\ZFC$ augmented by two inaccessible cardinals $\kappa_1 < \kappa_2$, providing Grothendieck universes $\mathbf{U}_1 \subset \mathbf{U}_2$. Sufficiency of this base for each proposed construction requires its own size and existence checks. The categorical machinery follows Kelly~\cite{Kelly1982} for enriched and $2$-categories, Lurie~\cite{LurieHTT, LurieHA} for $(\infty, 1)$-categories, Riehl--Verity~\cite{RiehlVerity2022} for $\infty$-cosmoi, Ad\'amek--Rosick\'y~\cite{AdamekRosicky1994} for accessible functors, and Barwick--Schommer-Pries~\cite{BarwickSchommerPries2021} for $(\infty, n)$-unicity. All $2$-categories are locally small and $\kappa_1$-accessible unless stated otherwise.

\section{\texorpdfstring{Rich-Metatheories and the Class $\cS_S$}{Rich-Metatheories and the Class S-S}}\label{sec:rs}

\begin{definition}[Rich-metatheory]\label{def:rs}
A first-order theory $S$ is a \textbf{Rich-metatheory} (R-S) if:
\begin{itemize}
\item \textbf{(R1)} $S$ interprets Robinson arithmetic $\mathsf{Q}$.
\item \textbf{(R2)} The set of axioms of $S$ is recursively enumerable over a countable language.
\item \textbf{(R3)} $S$ has a model $\cM$ in the chosen outer universe $\mathbf U_2$.
\item \textbf{(R4)} There is a recursive injection $\ulcorner \cdot \urcorner : \mathrm{Formulas}(S) \to \mathbb{N}$ such that the predicates ``$n$ codes a provable formula'' and ``$n$ codes a well-formed formula'' are $S$-definable.
\item \textbf{(R5)} There is $n_S \in \mathbb{N} \cup \{\infty\}$ such that the syntactic category $\Syn(S)$ (contexts, provable-from-context substitutions, higher homotopies) is an $(\infty, n_S)$-category, with the evaluation functor $\ev : \Syn(S) \to \cM$ preserving the $(\infty, n_S)$-structure up to $(\infty, n_S)$-equivalence.
\end{itemize}
We denote the class of Rich-metatheories as $\RS$, and the class of Rich-foundations (axiomatic systems satisfying (R1)--(R5)) as $\cF$.
\end{definition}

Examples must be checked against (R1)--(R5) using a specific formal theory and a chosen categorical semantics. A mathematical programme is not itself a recursive axiomatization, and an evaluation functor is not automatically fully faithful. The general higher-level syntax--semantics bridge is a hypothesis to verify in each application.

\begin{definition}[The selected presentation class $\cS_S$]\label{def:SS}
Fix a syntactic presentation $\Syn(S)$ and the allowed universe bounds. Include in $\cS_S^{\mathrm{local}}$ the pointed presentations $(\Syn(S),[\phi])$ for admissible formulas $\phi$ of its syntax, together with the model objects specified by the chosen interface. Define $\cS_S^{\mathrm{global}}$ using a stated list of actually existing, $S$-definable derived constructions. Set $\cS_S=\cS_S^{\mathrm{local}}\cup\cS_S^{\mathrm{global}}$. Existence, size, and closure under further limits are not consequences of this notation.
\end{definition}

\begin{lemma}[Membership of the chosen syntactic presentation]\label{lem:SS-membership}
If a formally described candidate is represented by the pointed presentation $(\Syn(S),[\phi_X])$ under this convention, that presentation belongs to $\cS_S$.
\end{lemma}
\begin{proof}
It is included in $\cS_S^{\mathrm{local}}$ by Definition~\ref{def:SS}.
\end{proof}
This is a statement about the selected formal presentation. Identifying it with a semantic object in a particular model requires an additional comparison, not merely interpretation of the formula.

\begin{definition}[Morita-reducibility]\label{def:morita}
An $(\infty, n)$-functor $F : X \to Y$ is an \emph{$(\infty, n)$-equivalence onto its image} if $F$ is fully faithful and essentially surjective onto the full sub-$(\infty, n)$-category of $Y$ given by $\image(F)$. We say $X$ is \emph{$(\infty, n)$-Morita-reducible} to $Y$, written $X \to_{M, n} Y$, if such a functor exists.
\end{definition}

\section{The Absolute Foundation No-Go Theorem}\label{sec:afnt}

\begin{theorem}[Relative AFN-T]\label{thm:afnt}
Fix a categorical setting and a selected class $\cS_S$. If a candidate $X$ belongs to $\cS_S$, it is Morita-reducible to an element of that class. Therefore it cannot satisfy
\[
 (\Pi_4):\quad X\not\to_{M,n}Y\quad\text{for every }Y\in\cS_S.
\]
In particular, the conjunction of such represented definability, $\Pi_4$, and maximal generativity is empty.
\end{theorem}
\begin{proof}
Take $Y=X$. The identity functor is an equivalence and hence a reduction prohibited by $\Pi_4$.
\end{proof}
The former proof incorrectly inferred full faithfulness of an arbitrary evaluation functor from an initiality property of syntax. Interpretation can identify distinct syntactic data; full faithfulness needs an independent theorem. The relative identity proof avoids that inference and states its actual scope.

\begin{theorem}[Relative boundary at a constructed limit]\label{thm:afnt-beta}
Suppose a specified $S$-definable small diagram has a colimit $A_\infty$ in the chosen universe and the selected derived class is closed under this construction. Then $A_\infty\in\cS_S$ and it fails $\Pi_4$.
\end{theorem}
\begin{proof}
Membership is the stated closure hypothesis; apply Theorem~\ref{thm:afnt}.
\end{proof}
A bound on cofinality alone does not construct a definable diagram or its colimit. A uniformly described proper-class scheme can have a finite description; it is not excluded merely by counting formula codes. Its semantics and consistency strength require a separate analysis.

\section{Absoluteness}\label{sec:absolute}

\begin{theorem}[Transfer of the relative obstruction]\label{thm:absolute}
At every stage or presentation where the candidate lies in the chosen comparison class and identity maps are admissible reductions, the conclusion of Theorem~\ref{thm:afnt} holds.
\end{theorem}
\begin{proof}
Apply the identity argument in that presentation.
\end{proof}
Transferring a claim between different metatheories, depths, or semantic models requires a translation preserving these hypotheses. This does not establish exhaustiveness of possible variations.

\begin{remark}[Unbounded stabilization remains open]\label{thm:stab}
Barwick--Schommer-Pries~\cite{BarwickSchommerPries2021} compares models of $(\infty,n)$-categories for fixed finite $n$. It does not prove an equivalence between unspecified $(\infty,\infty)$ and $(\infty,\infty+1)$ constructions. An unbounded-depth model, coherent comparison maps, and a proof of any proposed stabilization must be supplied separately.
\end{remark}

\section{Closure of Bypass Paths}\label{sec:bypass}

The following statements specify what can be concluded for constructed universe families, reflection sequences, and intensional data. They do not exclude every possible extension of these kinds.

\subsection{Universe Polymorphism}

\begin{theorem}[Definable coherent universe families]\label{thm:universe}
Let $D:I\to\Cat$ be a specified small $S$-definable diagram. If its relevant limit exists within the chosen universe and is included among the allowed derived constructions, that limit lies in $\cS_S$ and fails $\Pi_4$.
\end{theorem}
\begin{proof}
Use the closure hypothesis and Theorem~\ref{thm:afnt}.
\end{proof}
An endofunctor of a larger universe need not preserve a smaller one. Thus restriction does not automatically define transition functors between universe-endofunctor categories. Their preservation and coherence must be proved. An individual compatible family is also not identical to the entire category of compatible families.

\subsection{Reflective Towers}

\begin{remark}[Reflection requires a separate strength analysis]\label{thm:reflective}
The rule $S_{\alpha+1}=S_\alpha+\Con(S_\alpha)$ does not specify limit stages, an effective ordinal-notation system, or a uniform reflection principle. No general statement that reflection adds exactly one inaccessible or stabilizes at an expression $\Con(S)+\kappa$ is established. At each actually constructed stage, the relative identity obstruction remains available with that stage's own derived class; it does not compare proof-theoretic strengths or show termination.
\end{remark}

\subsection{Intensional Refinement}

\begin{definition}[Display data and the ordinary target]\label{def:Sint}
For the verified ordinary construction, a display category is a category with pullbacks and a selected class of arrows containing isomorphisms and closed under pullback, composition, and isomorphism of arrows. Let $\Sint$ be their $2$-category with display- and pullback-preserving functors and natural transformations. A higher target requires a separately specified model of arrows, transformations, and coherence.
\end{definition}

\begin{theorem}[Ordinary slice construction]\label{thm:intensional}
For an ordinary locally small category $\mathcal B$ with chosen pullbacks, the assignment $F\mapsto(\mathcal B/F,\text{isomorphisms})$, with base-change functors, defines a contravariant pseudofunctor $\mathcal B^{\mathrm{op}}\to\Sint$. Its display filtration generated from isomorphisms by the stated operations is constant at stage zero.
\end{theorem}
\begin{proof}
Slices have pullbacks and base change preserves them. Their universal property gives coherent identity and composition isomorphisms for reindexing. Every generating operation preserves isomorphisms, so it cannot add a nonisomorphism to the seed class.
\end{proof}
The full MSFS refinement still requires: (I-1) coherent action on higher cells; (I-2) a verified action of gauge transformations; (I-3) actual Morita-equivalent witnesses separated by retained intensional data; and (I-4) a correctly typed comparison of matching variance and a proved localization equivalence.

The former bare-$\MLTT$/$\ETT$ witness is insufficient. The relevant type-theoretic comparison requires uniqueness of identity proofs and functional extensionality~\cite{KapulkinLi2025}. Normalization, typing, and conversion also require specific effective presentations to be transported. Yoneda alone does not identify $\cF/F$ with $\rho(F)$.

Given a comparison $\widetilde\pi I\simeq\pi$ with compatible domains, the base image of $I(F)$ is $\pi(F)$. A fibre must include its comparison-to-the-base data; nontriviality and a fibration structure are additional claims. Such base-image constancy does not prove absence of new computational or resource-sensitive structure.

\paragraph{Finite observation and search.}
An executable observational quotient fixes the interpreter, complete library
meanings, all admission evidence, and the continuation policy. Exact register
traces with a separate undefined outcome support a continuation congruence
only with all policy metadata, including constant-use and selection flags,
and equality-preserving updates. This preserves the specified finite checks,
not behavior on every new input or physical cost. Distinct representations
may remain distinct explanations; checking discovered final candidates before
merging prefixes does not enumerate the descendants of discarded alternatives.
Budget exhaustion gives \texttt{Unknown}, rather than semantic undefinedness
or refutation. These distinctions leave the higher intensional construction
obligations above open.

\section{Meta-Classification of Level 5+ Frameworks}\label{sec:meta}

\begin{definition}[Level 5+ meta-structure]\label{def:meta}
A $2$-category $\mathbf{F}$ is a \emph{Level 5+ meta-structure} if:
\begin{itemize}
\item \textbf{(M1)} locally small;
\item \textbf{(M2)} equipped with a classification $2$-functor $\Cl_{\mathbf{F}} : \mathbf{F} \to \fM(\mathbf{F}) \subseteq \fM$;
\item \textbf{(M3)} accessible reflection $R_{\mathbf{F}} : \mathbf{F} \to \mathbf{F}$;
\item \textbf{(M4)} non-generative: $\mathbf{F}$ does not derive the axioms of $F \in \image(\Cl_{\mathbf{F}})$;
\item \textbf{(M5)} parametrized by some $S \in \RS$.
\end{itemize}
$\mathbf{F}$ is \emph{maximal} (and we write $\mathbf{F} \in \Meta_{5+}^{\max}$) if additionally:
\begin{itemize}
\item \textbf{(Max-1)} $\image(\Cl_{\mathbf{F}}) = \fM$;
\item \textbf{(Max-2)} $\mathrm{Aut}_2(\mathbf{F})$ acts transitively on gauge classes in $\image(\Cl_{\mathbf{F}})$;
\item \textbf{(Max-3)} $\mathbf{F}$ admits a depth-indexed filtration $\{\mathbf{F}^{(m)}\}_{m < \omega}$ with $\mathbf{F} \simeq_2 \colim_m \mathbf{F}^{(m)}$ and a depth function $\mathrm{depth} : \Ob(\mathbf{F}) \to \omega$ such that every comprehension-defined object of depth $m$ uses only objects of depth $< m$;
\item \textbf{(Max-4)} a correctly typed higher intensional construction satisfying the additional (I-1)--(I-4) obligations after Theorem~\ref{thm:intensional}, with the joint-separation condition of the main paper. Its existence is not supplied by the ordinary slice construction.
\end{itemize}
The class of Level 5+ meta-structures is $\Meta_{5+}$; the class of maximal meta-structures is $\Meta_{5+}^{\max} \subseteq \Meta_{5+}$.
\end{definition}

\begin{theorem}[Conditional comparison criterion]\label{thm:meta-cat}
Suppose $2$-functors $\Phi_i:\mathbf F_i\to\mathbf K$, $i=1,2$, are essentially surjective on objects and induce equivalences on all hom-categories. Then $\mathbf F_1\simeq_2\mathbf F_2$.
\end{theorem}
\begin{proof}
Both comparisons are biequivalences. Compose one with a pseudoinverse of the other.
\end{proof}
The maximality axioms do not by themselves supply these comparisons or lift arbitrary morphisms. For example, $1\to BG$ for nontrivial $G$ is essentially surjective and faithful but not full. Models of an accessible theory need not all be equivalent; the theory of groups already has inequivalent models. A higher-depth equivalence requires further coherence data.

\begin{theorem}[Conditional image separation]\label{thm:meta-mult}
Two constructed classifiers over the same base with distinct essential classification images cannot be equivalent over that base. A classifier with a proper image fails (Max-1).
\end{theorem}
\begin{proof}
An equivalence commuting with the classification functors preserves their essential images.
\end{proof}
Applying this criterion to Univalent Foundations, $\infty$-cosmoi, or cohesive frameworks requires actual classifiers satisfying (M1)--(M5) and proofs of the proposed image distinctions. Programme names and different syntactic presentations do not furnish these witnesses.

\begin{remark}[Meta-classification remains a construction problem]\label{thm:meta-stab}
The proposed quotient $\Meta_{5+}/\mathrm{gauge}_{\mathrm{meta}}$ needs a site, descent data, size bounds, a classification functor of the correct type, and a reflection. Idempotence under further classification requires explicit comparison functors and invertible unit/counit. Neither fixed-finite-$n$ unicity nor the notation for the quotient supplies them. Two inaccessible cardinals also do not give an unbounded tower of inaccessible universes. At every stage that is actually constructed in its selected derived class, Theorem~\ref{thm:afnt} still applies independently of stabilization.
\end{remark}

\section{Consistency Strength}\label{sec:cons}

\begin{theorem}[Conditional definitional conservativity]\label{thm:consistency}
If all new apparatus is introduced by explicit definitions whose existence and uniqueness are proved in a fixed base $B$, adjoining its names is conservative over $B$ for sentences of the old language.
\end{theorem}
\begin{proof}
Eliminate the defined names throughout any proof, using existence and uniqueness for functional definitions.
\end{proof}
The proposed classifier and higher intensional constructions have not all been completed in $\ZFC+2\text{-inacc}$. No exact consistency-strength equation is therefore inferred here. A complete model construction would give an upper bound; equiconsistency additionally requires a reverse interpretation or a separate lower-bound proof.

\begin{corollary}[Level-6 impossibility]\label{cor:level-6}
Within the represented-definability convention of Definition~\ref{def:SS} and Lemma~\ref{lem:SS-membership}, no candidate satisfies that condition together with $\Pi_4$ and maximal generativity. This is the relative exclusion of Theorem~\ref{thm:afnt}.
\end{corollary}

\begin{proof}
Immediate from Theorem~\ref{thm:afnt}.
\end{proof}

\section*{Acknowledgments}

The author thanks the mathematical community for the foundational results of Cantor, Russell, G\"odel, Tarski, Lawvere, Lambek--Scott, Ad\'amek--Rosick\'y, Makkai--Par\'e, Jacobs, Pronk, Lurie, Riehl--Verity, Voevodsky, Schreiber, and Barwick--Schommer-Pries, on which this synthesis rests.

\begin{thebibliography}{99}

\bibitem{AdamekRosicky1994}
J.~Ad\'amek and J.~Rosick\'y,
\newblock \emph{Locally Presentable and Accessible Categories},
\newblock London Math. Soc. Lecture Notes 189, Cambridge Univ. Press, 1994.

\bibitem{AraMaltsiniotis2017}
D.~Ara and G.~Maltsiniotis,
\newblock Homotopy comparison of cubical and globular models of weak $\omega$-categories,
\newblock preprint, 2017.

\bibitem{AwodeyWarren2009}
S.~Awodey and M.~Warren,
\newblock Homotopy theoretic models of identity types,
\newblock \emph{Math. Proc. Cambridge Philos. Soc.} 146 (2009), 45--55.

\bibitem{Bakovic2010}
I.~Bakovi\'c,
\newblock \emph{Fibrations of bicategories},
\newblock preprint, 2010.

\bibitem{BarwickSchommerPries2021}
C.~Barwick and C.~Schommer-Pries,
\newblock On the unicity of the theory of higher categories,
\newblock \emph{J. Amer. Math. Soc.} 34 (2021), 1011--1058; arXiv:1112.0040.

\bibitem{Cantor1891}
G.~Cantor,
\newblock \"Uber eine elementare Frage der Mannigfaltigkeitslehre,
\newblock \emph{Jahresbericht der Deutschen Math.-Vereinigung} 1 (1891), 75--78.

\bibitem{ClausenScholze2019}
D.~Clausen and P.~Scholze,
\newblock \emph{Lectures on Condensed Mathematics},
\newblock Univ. of Bonn, 2019,
\newblock \url{https://www.math.uni-bonn.de/people/scholze/Condensed.pdf}.

\bibitem{Connes1994}
A.~Connes,
\newblock \emph{Noncommutative Geometry},
\newblock Academic Press, 1994.

\bibitem{Ernst2015}
M.~Ernst,
\newblock The prospects of unlimited category theory: doing what remains to be done,
\newblock \emph{Rev. Symb. Log.} 8 (2015), 306--327.

\bibitem{Feferman1964}
S.~Feferman,
\newblock Systems of predicative analysis,
\newblock \emph{J. Symb. Log.} 29 (1964), 1--30.

\bibitem{Godel1931}
K.~G\"odel,
\newblock \"Uber formal unentscheidbare S\"atze der Principia Mathematica und verwandter Systeme I,
\newblock \emph{Monatsh. Math. Phys.} 38 (1931), 173--198.

\bibitem{Hermida1999}
C.~Hermida,
\newblock Some properties of Fib as a fibred $2$-category,
\newblock \emph{J. Pure Appl. Algebra} 134 (1999), 83--109.

\bibitem{KapulkinLi2025}
C.~Kapulkin and Y.~Li,
\newblock Extensional concepts in intensional type theory, revisited,
\newblock \emph{Theoretical Computer Science} 1029 (2025), 115051; arXiv:2310.05706.

\bibitem{Hofmann1995}
M.~Hofmann,
\newblock \emph{Extensional concepts in intensional type theory},
\newblock Ph.D. thesis, Univ. of Edinburgh, 1995.

\bibitem{HoTTBook2013}
The Univalent Foundations Program,
\newblock \emph{Homotopy Type Theory: Univalent Foundations of Mathematics},
\newblock Institute for Advanced Study, 2013,
\newblock \url{https://homotopytypetheory.org/book/}.

\bibitem{Hyland1982}
J.M.E.~Hyland,
\newblock The effective topos,
\newblock in \emph{The L.E.J. Brouwer Centenary Symposium}, North-Holland, 1982, pp.~165--216.

\bibitem{KapulkinLumsdaine2021}
K.~Kapulkin and P.L.~Lumsdaine,
\newblock The simplicial model of univalent foundations (after Voevodsky),
\newblock \emph{J. Eur. Math. Soc.} 23 (2021), 2071--2126; arXiv:1211.2851.

\bibitem{Kelly1982}
G.M.~Kelly,
\newblock \emph{Basic Concepts of Enriched Category Theory},
\newblock London Math. Soc. Lecture Notes 64, Cambridge Univ. Press, 1982.

\bibitem{Lair1981}
C.~Lair,
\newblock Cat\'egories modelables et cat\'egories esquissables,
\newblock \emph{Diagrammes} 6 (1981), L1--L20.

\bibitem{Lawvere1969}
F.W.~Lawvere,
\newblock Adjointness in foundations,
\newblock \emph{Dialectica} 23 (1969), 281--296.

\bibitem{LurieHTT}
J.~Lurie,
\newblock \emph{Higher Topos Theory},
\newblock Annals of Math. Studies 170, Princeton Univ. Press, 2009.

\bibitem{LurieHA}
J.~Lurie,
\newblock \emph{Higher Algebra},
\newblock preprint, 2017.

\bibitem{MakkaiPare1989}
M.~Makkai and R.~Par\'e,
\newblock \emph{Accessible Categories: The Foundations of Categorical Model Theory},
\newblock Contemp. Math. 104, Amer. Math. Soc., 1989.

\bibitem{Pronk1996}
D.~Pronk,
\newblock Etendues and stacks as bicategories of fractions,
\newblock \emph{Compos. Math.} 102 (1996), 243--303.

\bibitem{Rathjen1999}
M.~Rathjen,
\newblock The realm of ordinal analysis,
\newblock in \emph{Sets and Proofs}, Cambridge Univ. Press, 1999, pp.~219--279.

\bibitem{Rezk2010}
C.~Rezk,
\newblock A Cartesian presentation of weak $n$-categories,
\newblock \emph{Geom. Topol.} 14 (2010), 521--571; arXiv:0901.3602.

\bibitem{RiehlVerity2022}
E.~Riehl and D.~Verity,
\newblock \emph{Elements of $\infty$-Category Theory},
\newblock Cambridge Univ. Press, 2022.

\bibitem{Russell1903}
B.~Russell,
\newblock \emph{The Principles of Mathematics},
\newblock Cambridge Univ. Press, 1903.

\bibitem{Schreiber2013}
U.~Schreiber,
\newblock \emph{Differential cohomology in a cohesive $\infty$-topos},
\newblock arXiv:1310.7930, 2013.

\bibitem{Streicher1991}
T.~Streicher,
\newblock \emph{Semantics of type theory},
\newblock Birkh\"auser, 1991.

\bibitem{Tarski1936}
A.~Tarski,
\newblock Der Wahrheitsbegriff in den formalisierten Sprachen,
\newblock \emph{Studia Philos.} 1 (1936), 261--405.

\bibitem{Voevodsky2010}
V.~Voevodsky,
\newblock \emph{Foundations} Coq library,
\newblock software, 2010--2015,
\newblock \url{https://github.com/vladimirias/Foundations}.

\bibitem{Werner1994}
B.~Werner,
\newblock \emph{Une Th\'eorie des Constructions Inductives},
\newblock Ph.D. thesis, Univ. Paris 7, 1994.

\end{thebibliography}

\end{document}
